% =============================================================
%  BML TOOLBOX REFERENCE GUIDE  v1.0
%  Brain Modulation Laboratory · MGH / University of Pittsburgh
%  Generated: April 2026
% =============================================================
\documentclass[11pt, a4paper]{report}

\usepackage{fontspec}
\setmainfont{Georgia}
\setmonofont[Scale=0.85]{Menlo}

\usepackage[top=2.8cm, bottom=2.8cm, left=3.0cm, right=3.0cm]{geometry}
\usepackage{microtype}
\usepackage{parskip}
\setlength{\parskip}{6pt}

\usepackage[dvipsnames,table]{xcolor}
\definecolor{bmlblue}{RGB}{0,72,144}
\definecolor{boxbg}{RGB}{242,246,252}
\definecolor{codefg}{RGB}{30,30,50}
\definecolor{codebg}{RGB}{248,248,252}
\definecolor{rulecolor}{RGB}{180,200,230}
\definecolor{warnyellow}{RGB}{255,248,220}
\definecolor{warnborder}{RGB}{204,153,0}
\definecolor{greenbg}{RGB}{240,252,240}
\definecolor{greenborder}{RGB}{80,160,80}
\definecolor{redbg}{RGB}{255,245,245}
\definecolor{redborder}{RGB}{180,60,60}
\definecolor{glossbg}{RGB}{250,248,240}
\definecolor{partcolor}{RGB}{0,72,144}
\definecolor{annotA}{RGB}{30,100,200}
\definecolor{annotB}{RGB}{180,40,40}
\definecolor{annotC}{RGB}{40,140,60}
\definecolor{annotD}{RGB}{160,80,160}

\usepackage{amsmath, amssymb, mathtools}

\usepackage{tikz}
\usetikzlibrary{shapes.geometric,arrows.meta,positioning,fit,backgrounds,
                calc,decorations.pathreplacing,patterns}
\usepackage{graphicx}
\usepackage{float}

\usepackage{booktabs}
\usepackage{tabularx}
\usepackage{longtable}
\usepackage{array}
\newcolumntype{L}[1]{>{\raggedright\arraybackslash}p{#1}}
\newcolumntype{C}[1]{>{\centering\arraybackslash}p{#1}}

\usepackage{listings}
\lstdefinestyle{matlab}{
  language=Matlab,
  backgroundcolor=\color{codebg},
  basicstyle=\ttfamily\small\color{codefg},
  keywordstyle=\color{bmlblue}\bfseries,
  commentstyle=\color{gray}\itshape,
  stringstyle=\color{Maroon},
  numbers=left, numberstyle=\tiny\color{gray}, numbersep=8pt,
  frame=single, framerule=0.4pt, rulecolor=\color{rulecolor},
  breaklines=true, captionpos=b, showstringspaces=false, tabsize=2
}

\usepackage[most]{tcolorbox}
\newtcolorbox{notebox}{colback=greenbg,colframe=greenborder,
  fonttitle=\bfseries,title=Note,left=4pt,right=4pt,top=3pt,bottom=3pt}
\newtcolorbox{warnbox}{colback=warnyellow,colframe=warnborder,
  fonttitle=\bfseries,title=Warning,left=4pt,right=4pt,top=3pt,bottom=3pt}
\newtcolorbox{formulabox}[1]{colback=boxbg,colframe=bmlblue!50,
  fonttitle=\bfseries\color{bmlblue},title=#1,left=4pt,right=4pt,top=3pt,bottom=3pt}
\newtcolorbox{glossentry}[1]{colback=glossbg,colframe=bmlblue!30,
  fonttitle=\bfseries\color{bmlblue},title=#1,left=5pt,right=5pt,top=4pt,bottom=4pt,
  before upper={\setlength{\parskip}{3pt}}}

\usepackage{enumitem}

\usepackage{hyperref}
\hypersetup{colorlinks=true,linkcolor=bmlblue,urlcolor=bmlblue,
            pdftitle={BML Toolbox Reference Guide v1.0},
            pdfauthor={Brain Modulation Laboratory}}

% Chapter styling
\usepackage{titlesec}
\titleformat{\part}[display]
  {\centering\normalfont\Huge\bfseries\color{partcolor}}
  {\partname\ \thepart}{20pt}{}
\titleformat{\chapter}[hang]
  {\normalfont\Large\bfseries\color{bmlblue}}
  {\thechapter}{12pt}{}[\vspace{-4pt}\rule{\textwidth}{0.4pt}]
\titlespacing*{\chapter}{0pt}{20pt}{10pt}

\usepackage{fancyhdr}
\pagestyle{fancy}\fancyhf{}
\fancyhead[L]{\small\color{gray}BML Toolbox Reference Guide}
\fancyhead[R]{\small\color{gray}\leftmark}
\fancyfoot[C]{\small\thepage}
\renewcommand{\headrulewidth}{0.3pt}
\fancypagestyle{plain}{\fancyhf{}\fancyfoot[C]{\small\thepage}}

% Reusable TikZ styles
\tikzset{
  annot/.style={rectangle,minimum height=0.45cm,inner sep=0pt,
                draw=#1!70,fill=#1!20,font=\tiny},
  taxis/.style={->,thick,gray!60},
  tlabel/.style={font=\tiny,gray},
  resultbar/.style={rectangle,minimum height=0.45cm,inner sep=0pt,
                    draw=annotC!80,fill=annotC!25,font=\tiny},
  removedbar/.style={rectangle,minimum height=0.45cm,inner sep=0pt,
                     draw=gray!50,fill=gray!10,font=\tiny,text=gray!60},
}

% Version info
\newcommand{\bmlversion}{1.0}
\newcommand{\bmldate}{April 2026}

\begin{document}

% ---- Title page ----
\begin{titlepage}
\centering
\vspace*{2cm}
{\color{bmlblue}\rule{\textwidth}{2pt}}\\[0.8em]
{\fontsize{28}{34}\selectfont\bfseries\color{bmlblue} BML Toolbox}\\[0.4em]
{\fontsize{18}{22}\selectfont\color{gray} Reference Guide}\\[0.3em]
{\color{bmlblue}\rule{\textwidth}{2pt}}\\[1.5em]
{\Large\bfseries Version \bmlversion}\\[0.5em]
{\large\color{gray}\bmldate}\\[2em]

\begin{tcolorbox}[colback=boxbg,colframe=bmlblue!40,width=12cm]
\centering\small
\textbf{Brain Modulation Laboratory}\\
Massachusetts General Hospital · University of Pittsburgh\\[0.5em]
\textit{A comprehensive reference for annotation table manipulation\\
and multi-device clock synchronization in MATLAB.}
\end{tcolorbox}

\vfill
\begin{tabular}{ll}
\toprule
\textbf{Document} & BML Toolbox Reference Guide\\
\textbf{Version} & \bmlversion\\
\textbf{Date} & \bmldate\\
\textbf{Language} & MATLAB (Python migration in progress)\\
\textbf{Devices supported} & Ripple/Trellis, Zoom, NeuroOmega, Psychtoolbox\\
\bottomrule
\end{tabular}
\end{titlepage}

% ---- Changelog ----
\chapter*{Version History}
\addcontentsline{toc}{chapter}{Version History}
\begin{center}
\begin{tabular}{L{1.5cm} L{2.5cm} L{9cm}}
\toprule
\textbf{Version} & \textbf{Date} & \textbf{Changes} \\
\midrule
1.0 & April 2026 &
  First unified release. Merged annotation algebra guide (Part~I)
  with synchronization pipeline guide (Part~II).
  Added: digital-trigger audio sync, task-event fallback strategy,
  old vs.\ new DP matching comparison, Python migration roadmap,
  best practices, complete worked examples, ReadTheDocs site. \\
\midrule
0.2 & March 2026 &
  Sync guide: added Psychtoolbox sync, expanded consolidation section,
  data flow diagram. \\
0.1 & January 2026 &
  Initial sync guide draft covering Zoom analog sync,
  NeuroOmega event sync, bml\_timealign/bml\_timewarp internals. \\
\bottomrule
\end{tabular}
\end{center}

% ---- Introduction ----
\chapter*{Introduction}
\addcontentsline{toc}{chapter}{Introduction}

\section*{About this guide}

This guide is the complete reference for the \textbf{BML (Brain Modulation Laboratory) MATLAB toolbox}, used for intraoperative and clinical neurophysiology data analysis at MGH and the University of Pittsburgh. It covers:

\begin{description}[leftmargin=2cm,style=nextline]
  \item[Part~I: Annotation Tables] The universal data structure used throughout BML. Every time-stamped object --- a trial, a spike, a file, an artifact window --- is an annotation table. Part~I covers the schema, all set/interval algebra operations (filter, intersect, union, difference, consolidate, \ldots), I/O, visualization, and a complete worked example.
  \item[Part~II: Synchronization] How to align recordings from four independent devices (Ripple/Trellis, Zoom audio recorder, NeuroOmega, task laptop) onto a single master clock. Covers analog cross-correlation, digital event matching, the DP algorithm, the fallback strategy, consolidation, and best practices.
\end{description}

\section*{Design philosophy}

Three principles underlie every BML function:

\begin{enumerate}
  \item \textbf{Annotation tables are the universal contract.} Every function takes annotation tables as input and returns annotation tables as output. This makes functions composable: the output of \texttt{bml\_annot\_filter} feeds directly into \texttt{bml\_annot\_intersect} without any data transformation.
  \item \textbf{All times are on the master clock.} After synchronization, every sample index and every event timestamp is expressed in seconds on the Trellis/Ripple hardware clock. You never need to know which device produced a data point.
  \item \textbf{Non-destructive.} Raw files are never modified. Sync state is stored separately in \texttt{sync\_roi} tables. Any step can be re-run without corrupting the original data.
\end{enumerate}

\section*{How to use this guide}

\begin{itemize}
  \item \textbf{New to BML?} Start with Section~\ref{sec:schema} (annotation table schema), then Section~\ref{sec:roi} (the ROI table and coordinate system), then the worked example in Chapter~\ref{sec:example}.
  \item \textbf{Doing sync?} Jump to Part~II. The sync overview diagram (Section~\ref{sec:dataflow}) shows the full pipeline at a glance.
  \item \textbf{Looking for a specific function?} Use the reference table in Chapter~\ref{sec:reference} or the table of contents.
\end{itemize}

\tableofcontents
\newpage

% ================================================================
\part{Annotation Tables and Interval Algebra}
% ================================================================

\input{/tmp/annot_body}

% ================================================================
\part{The Synchronization Pipeline}
% ================================================================

\input{/tmp/sync_body}

\end{document}
